\section{Comparison with $C^*$-algebras}
\label{sec:cstar}

\subsection{The two embeddings of classical geometry}

Both theories extend the same commutative picture, along different embeddings:
\[
  F_{\Cstar}=C(-,\CC),\qquad F_{\PB}=C(-,\RR),
  \qquad\text{related by }\ \Lambda=(-)\otimes_\RR\CC,\quad \Lambda\circ F_{\PB}\cong F_{\Cstar}.
\]
Any comparison must be taken relative to $\Lambda$: as a \emph{real} algebra,
$C(X,\CC)$ is not a PB-algebra, since it contains the central element $i$ and
fails \eqref{ax:pb2} there (Example~\ref{ex:excluded}).

\begin{definition}\label{def:intersection}
Let $\mathcal{I}$ be the class of PB-algebras $A$ whose complexification
$A_\CC$ carries a compatible $C^*$-structure --- equivalently, the real forms of
$C^*$-algebras that satisfy \eqref{ax:pb1} and \eqref{ax:pb2}.
\end{definition}

$\mathcal{I}$ is populated: $M_2(\RR)$ and $\HH$ both complexify to $M_2(\CC)$,
and both are PB-algebras with constant $\tfrac14$ (Examples~\ref{ex:H},
\ref{ex:M2}).  Note that these are two different real forms of one
$C^*$-algebra, so $\Lambda$ is far from injective on objects.  The relationship
between the theories is a span
\[
  \PB\;\hookleftarrow\;\mathcal{I}\;\longrightarrow\;\Cstar,
\]
and not an inclusion in either direction --- with a caveat on the first half
that the earlier notes do not record.  Proposition~\ref{prop:chirality} gives
$C^*$-algebras outside the image (no real form at all), at every constant.  In
the other direction, Proposition~\ref{prop:T2} gives a PB-object outside
$\mathcal{I}$ --- $T_2(\RR)$ has non-central radical, impossible in a semisimple
$C^*$-algebra --- but only for $C\ge\kap^*$, since below that $T_2(\RR)$ is not
an object (Proposition~\ref{prop:T2-constant}).  Whether
$\PB^{1/4}_{\gauge_1}\subseteq\mathcal{I}$, i.e.\ whether the span degenerates to
an inclusion at the critical constant, is open and is implied by
Conjecture~\ref{conj:semisimple} only in part: semisimplicity is necessary for
membership in $\mathcal{I}$, not sufficient.

\subsection{Forgetting the involution}

\begin{proposition}\label{prop:U-faithful}
Let $U$ be the forgetful assignment taking a real form of a $C^*$-algebra, with
its $*$-homomorphisms, to the underlying PB-algebra with its PB-morphisms.
Then $U$ is faithful and bijective on objects but \emph{not full}: there are
PB-morphisms between objects of $\mathcal{I}$ that are not induced by
$*$-homomorphisms.
\end{proposition}

\begin{proof}[Proof sketch]
Faithfulness and bijectivity on objects are immediate.  For fullness, take
$A=M_2(\RR)$ and $g\in GL_2(\RR)$ with $\Ad_g$ isometric and gauge-preserving
but $g$ not a scalar multiple of an orthogonal matrix; then $\Ad_g$ is a
PB-automorphism of $A$ whose complexification is not a $*$-automorphism of
$M_2(\CC)$, since $*$-automorphisms of $M_2(\CC)$ are exactly the $\Ad_v$ with
$v$ unitary.
\end{proof}

\begin{remark}
So on the intersection the PB-theory has \emph{strictly more} morphisms off the
centre, and $\mathcal{I}$ carries two different subobject lattices:
$*$-subalgebras on one side, saturated subalgebras
(Definition~\ref{def:saturated}) on the other.  $U$ is an equivalence exactly on
the commutative floor, where the gauge vanishes identically and every unital
homomorphism is legal.  $\mathcal{I}$ is a span-vertex, not a subcategory of
either theory in a way that respects both structures --- and in particular there
is no reason for it to be closed under PB-quotients.
\end{remark}

\subsection{A uniform bound on the $C^*$ side}

The examples suggest something stronger than membership case by case.

\begin{conjecture}[Uniform bound]\label{conj:cstar-uniform}
Every real $C^*$-algebra $A$ --- a real form of a $C^*$-algebra with its
inherited norm --- satisfies \eqref{ax:pb1} with gauge $\gauge_1$ and constant
$C=\tfrac14$.  Equivalently, by Proposition~\ref{prop:real-stampfli},
\[
  \norm{a}^2-\norm{a^2}\;\le\;\dist(a,\RR1)^2\qquad\text{for all }a .
\]
\end{conjecture}

\begin{remark}
The evidence is as follows.  Equality holds at every square-zero element: if
$a^2=0$ and $\norm{a}=1$ then the left side is $1$, and writing $a$ in its
block form shows $\norm{a-t1}\ge\norm{a}$ for every $t$, so the right side is
$1$ as well.  The inequality also holds at every normal element, where the left
side vanishes.  Numerically it survives a search over $M_n(\RR)$ for
$n\le4$ (\S\ref{sec:status}).  A proof would be valuable well beyond this
subject: it says that failure of the square property is controlled by distance
to the scalars, uniformly, with the sharp constant attained on nilpotents.
Note that \eqref{ax:pb2} is a genuinely separate matter and is \emph{not} implied:
$A=\CC$ is a real $C^*$-algebra failing \eqref{ax:pb2}, and so is $M_2(\CC)$
viewed as a real algebra, whose centre $\CC$ is not of the form $C(X,\RR)$.  The
gate is always the centre.
\end{remark}
